% -*- mode : latex; coding: utf-8 -*-
\chapter{Introduction}

Modern datacenter workloads like deep learning (DL), artificial intelligence (AI) have increasing demand for large memory capacity. However, modern processors have limited number of pins to add DIMM (Dual in-line memory module). The limited number of pins leads to imbalance between peak computation and memory capacity. If this problem isn't addressed, future datacenter performance will be limited by inadequate memory capacity.

In traditional memory design, memory modules are co-located on a system board, restricting the configuration and scalability of both compute and memory resources \cite{disaggregate}. In this design, it's not possible to add more memory independently. To overcome this challenge, disaggregated memory has been suggested. Disaggregated memory design utilizes shared memory blades that can be accessed, as needed, by multiple compute blades via shared blade interconnect \cite{disaggregate}. This design enables independent scaling of memory capacity or compute power. Thus, datacenter processors can overcome inadequate memory problem.

Disaggregated memory requires interconnection between processors and disaggregated memory. The emerging Compute Express Link (CXL) interconnect standard \cite{CXL} greatly facilitates fast and deployable disaggregated memory \cite{first-generation}. CXL enables native load/store (ld/st) accesses to disaggregated memory for commercial processors and memory devices. CXL has become one of the best option for memory disaggregation by superceding competing protocols \cite{first-generation}.

In this research, I implemented cycle-accurate and trace-driven CXL simulator \textit{CXLmem}. This simulator emulates CXL interconnection, especially CXL type 3 device that exploits CXL.mem protocol. By using this simulator, I evaluated the performance of CXL attached memory with general benchmarks such as SPEC CPU 2006. This research will contribute to further research about memory disaggregation via CXL interconnection.

Here are the key contributions of this research.
\begin{itemize}
    \item CXLmem: I implemented the cycle-accurate and trace-driven CXL type 3 device simulator \textit{CXLmem}. This simulator helps understanding the detailed performance of CXL attached memory. This simulator can be simulated as standalone simulator or as memory module attached to system architecture simulator such as MacSim \cite{MACSIM}. When it runs as standalone simulator, if we feed memory trace to the simulator, we can get a detailed performance of CXL attached memory. On the other hand, if we attach it to system architecture simulator, we can see the impact of CXL attached memory to overall system. In this case, we can also get detailed performance of CXL attached memory. By adding some more features to CXLmem, we can develop CXL type 2 device simulator and use it for futher research.
    \item Performance of CXL attached memory: I evaluated the performance of CXL attached memory by using CXLmem. CXLmem provides memory related performance and CXL interconnection related performance such as memory read latency. This evaluation results will help understanding the impact of CXL attached memory to system and studying CXL attached memory in the future work.
\end{itemize}

In this paper, I first provide background about CXL interconnection at Chapter \ref{sec:backgrounds}. After that I explain the architecture and implementation of CXL type 3 device simulator CXLmem in Chapter \ref{sec:CXLmem Design}. In Chapter \ref{sec:evaluation}, I show the CXL type 3 device performance such as IPC (instruction per cycle), read latency evaluated using CXLmem. Finally, Chapter \ref{sec:conclusion} concludes and suggests possible future research directions.